\documentclass[11pt,a4paper]{article}
\usepackage[margin=1in]{geometry}
\usepackage{amsmath,amssymb,amsthm}
\usepackage{algorithm}
\usepackage{algpseudocode}
\usepackage{booktabs}
\usepackage{hyperref}

\theoremstyle{definition}
\newtheorem{definition}{Definition}[section]
\newtheorem{theorem}{Theorem}[section]
\newtheorem{lemma}[theorem]{Lemma}
\newtheorem{remark}{Remark}[section]

\title{\textbf{Theoretical Foundations, Channel Recalibration, and Spatial Gating of Squeeze-and-Excitation and CBAM}}
\author{Team Scaibu \\ \small Email: \texttt{jaydeep.9664920749@gmail.com} \\ \small Architecture and Core Engine Team}
\date{October 2026}

\begin{document}
\maketitle

\begin{abstract}
This document presents the formal mathematical specification, global contextual pooling dynamics, and multi-dimensional attention gating of Squeeze-and-Excitation (SE) Networks and the Convolutional Block Attention Module (CBAM). Convolutional layers inherently treat all feature channels and spatial regions uniformly. Squeeze-and-Excitation introduces channel-wise attention by aggregating global spatial context via global average pooling and computing non-linear channel recalibration weights through a bottleneck MLP. CBAM extends this paradigm by sequentially applying dual-pooled channel attention followed by spatial attention gating. We formulate the attention operators, prove representational capacity bounds, analyze operational complexities, trace a worked numerical example, and document integration practices.
\end{abstract}

\section{Notation and Mathematical Preliminaries}

Let $X \in \mathbb{R}^{C \times H \times W}$ denote the input feature activation tensor.

\begin{itemize}
    \item $C \in \mathbb{N}_{\ge 1}$: Number of feature channels.
    \item $H, W \in \mathbb{N}_{\ge 1}$: Spatial height and width dimensions.
    \item $r \in \mathbb{N}_{\ge 1}$: Channel reduction ratio ($r = 16$ default).
    \item $C_{\text{mid}} = \lfloor C / r \rfloor$: Bottleneck hidden dimension in attention MLP.
    \item $W_1 \in \mathbb{R}^{C_{\text{mid}} \times C}, W_2 \in \mathbb{R}^{C \times C_{\text{mid}}}$: Excitation MLP projection matrices.
    \item $\sigma(z) = \frac{1}{1 + e^{-z}}$: Sigmoid activation function bounding gating scalars into $[0, 1]$.
    \item $M_c \in [0, 1]^C$: 1D channel attention scaling vector.
    \item $M_s \in [0, 1]^{H \times W}$: 2D spatial attention gating map.
    \item $Y \in \mathbb{R}^{C \times H \times W}$: Recalibrated output tensor.
\end{itemize}

\begin{table}[h]
\centering
\caption{Summary of Attention Mechanics in SE and CBAM}
\begin{tabular}{lll}
\toprule
\textbf{Module} & \textbf{Spatial Reduction} & \textbf{Attention Formulation} \\
\midrule
\textbf{SE} & Global Average Pooling & $M_c = \sigma(W_2 \cdot \text{ReLU}(W_1 \cdot \text{GAP}(X)))$ \\
\textbf{CBAM Channel} & AvgPool + MaxPool & $M_c = \sigma(\text{MLP}(\text{GAP}(X)) + \text{MLP}(\text{GMP}(X)))$ \\
\textbf{CBAM Spatial} & Channel Avg + Max & $M_s = \sigma(f^{7 \times 7}([\text{AvgPool}_c(X') \mathbin{\Vert} \text{MaxPool}_c(X')]))$ \\
\bottomrule
\end{tabular}
\end{table}

\section{Problem Definition and Core Concepts}

\subsection{Intuitive Meaning}
Standard convolutions operate within fixed local receptive fields ($3 \times 3$) and cannot exploit contextual dependencies across distant channels. Squeeze-and-Excitation allows the network to adaptively boost informative feature maps (e.g. object category indicators) while suppressing redundant background channels. CBAM tells the network "what" to pay attention to along channels, and "where" to focus spatially.

\subsection{Formal Mathematical Formulations}

\subsubsection{Squeeze-and-Excitation (SE)}
\begin{align}
z_c &= \frac{1}{H W} \sum_{h=1}^H \sum_{w=1}^W X(c, h, w) \quad (\text{Squeeze}) \\
s &= \sigma\left( W_2 \cdot \max(0, W_1 \cdot z) \right) \quad (\text{Excitation}) \\
Y(c, h, w) &= s_c \cdot X(c, h, w) \quad (\text{Recalibration})
\end{align}

\subsubsection{CBAM Attention Sequential Pipeline}
\begin{align}
M_c(X) &= \sigma\left( W_2 \text{ReLU}(W_1 \text{AvgPool}(X)) + W_2 \text{ReLU}(W_1 \text{MaxPool}(X)) \right) \\
X' &= M_c(X) \odot X \\
M_s(X') &= \sigma\left( f^{7 \times 7}\left( \left[ \frac{1}{C}\sum_{c=1}^C X'_c \mathbin{\Big\Vert} \max_c X'_c \right] \right) \right) \\
Y &= M_s(X') \odot X'
\end{align}

\section{Algorithm Specification}

\begin{algorithm}[H]
\caption{Squeeze-and-Excitation and CBAM Recalibration Pass}
\label{alg:se_cbam}
\begin{algorithmic}[1]
\Require Input $X \in \mathbb{R}^{C \times H \times W}$, MLP weights $W_1, W_2$, mode $M \in \{\text{se}, \text{cbam}\}$, spatial weights $W_s$.
\Ensure Output tensor $Y \in \mathbb{R}^{C \times H \times W}$, channel weights $s$, spatial weights $M_s$.
\If{$M = \text{se}$}
    \State $z \gets \text{GlobalAveragePool}(X)$
    \State $h \gets \text{ReLU}(W_1 z)$
    \State $s \gets \text{Sigmoid}(W_2 h)$
    \State $Y \gets s \odot X$
    \State \Return $\{Y, s\}$
\Else
    \State $z_{\text{avg}} \gets \text{GlobalAveragePool}(X)$, $z_{\text{max}} \gets \text{GlobalMaxPool}(X)$
    \State $s \gets \text{Sigmoid}(\text{MLP}(z_{\text{avg}}) + \text{MLP}(z_{\text{max}}))$
    \State $X' \gets s \odot X$
    \State $F_{\text{avg}} \gets \text{MeanChannel}(X')$, $F_{\text{max}} \gets \text{MaxChannel}(X')$
    \State $M_s \gets \text{Sigmoid}(\text{Conv2D}([F_{\text{avg}} \mathbin{\Vert} F_{\text{max}}], W_s))$
    \State $Y \gets M_s \odot X'$
    \State \Return $\{Y, s, M_s\}$
\EndIf
\end{algorithmic}
\end{algorithm}

\section{Theoretical Guarantees and Representational Capacity}

\begin{theorem}[Non-Linear Channel Selectivity and Bound]
For any input $X$, the SE channel scaling vector $s \in (0, 1)^C$ acts as a bounded soft gating mechanism. The gradient of the recalibrated output with respect to input $X_{c, h, w}$ satisfies:
\begin{equation}
\frac{\partial Y_{c, h, w}}{\partial X_{c, h, w}} = s_c + X_{c, h, w} \frac{\partial s_c}{\partial X_{c, h, w}}
\end{equation}
enabling self-attention dynamics with less than $1\%$ parameter increase over baseline ResNet backbones.
\end{theorem}

\begin{proof}
Directly differentiating $Y_{c, h, w} = s_c X_{c, h, w}$ by product rule.
The bottleneck MLP requires only $2 \frac{C^2}{r}$ parameters. For $C = 256, r = 16$, parameters $= 2(256)(16) = 8192$ floats, negligible compared to a $3 \times 3$ conv layer requiring $256 \times 256 \times 9 = 589,824$ parameters ($1.38\%$).
\end{proof}

\subsection{Limitations}
\begin{itemize}
    \item \textbf{Global Spatial Averaging Blindness}: Averaging entire spatial dimensions into a scalar loses relative spatial positioning cues, which is why CBAM incorporates subsequent spatial attention.
\end{itemize}

\section{Complexity Analysis}

\subsection{Time Complexity}
\begin{itemize}
    \item \textbf{Total Time Complexity}: $\Theta(C \cdot H \cdot W + \frac{2}{r} C^2)$ FLOPs.
\end{itemize}

\subsection{Space Complexity}
\begin{itemize}
    \item \textbf{Storage}: $\mathcal{O}(C \cdot H \cdot W)$ memory for output feature tensor.
\end{itemize}

\section{Exactness and Error Analysis}

Reductions, MLP projections, and sigmoid gating are mathematically exact.

\section{Worked Numerical Example}

Let $C = 2, H = 1, W = 1, C_{\text{mid}} = 1$.
Input $X = [[[2.0]], [[4.0]]]$.
Weights: $W_1 = [[0.5, 0.5]], W_2 = [[1.0], [1.0]]$.

\begin{enumerate}
    \item \textbf{GAP}: $z = [2.0, 4.0]^T$.
    \item \textbf{MLP Layer 1 + ReLU}: $h = \max(0, 0.5(2) + 0.5(4)) = \max(0, 3.0) = 3.0$.
    \item \textbf{MLP Layer 2}: $l = [1(3), 1(3)] = [3.0, 3.0]$.
    \item \textbf{Sigmoid}: $s = [\sigma(3), \sigma(3)] = [0.952574, 0.952574]$.
    \item \textbf{Recalibration}:
    $Y_0 = 0.952574(2.0) \approx 1.905148, Y_1 = 0.952574(4.0) \approx 3.810296$.
\end{enumerate}

\section{Numerical Considerations and Edge Cases}

\begin{itemize}
    \item \textbf{Sigmoid Saturation}: Extreme logits ($|z| > 60$) are clamped in the reference implementation to prevent floating-point underflow/overflow.
\end{itemize}

\begin{thebibliography}{9}
\bibitem{hu2018squeeze} J.~Hu, L.~Shen, and G.~Sun, ``Squeeze-and-Excitation Networks,'' in \emph{IEEE Conference on Computer Vision and Pattern Recognition (CVPR)}, pp.~7132--7141, 2018.
\bibitem{woo2018cbam} S.~Woo, J.~Park, J.~Y.~Lee, and I.~S.~Kweon, ``CBAM: Convolutional Block Attention Module,'' in \emph{European Conference on Computer Vision (ECCV)}, pp.~3--19, 2018.
\end{thebibliography}

\end{document}
